\section{Introduction}\label{sec: introduction}
\subsection{Results}\label{subsec: intro results}
This work concerns the construction of a six-functor formalism for solid quasicoherent sheaves in spectral algebraic geometry. The geometric objects of this setting, the spectral algebraic stacks, carry a natural theory of quasicoherent sheaves and their cohomology. Our main result, stated here with some imprecision, is that once quasicoherent sheaves are replaced by their solid analogues from condensed mathematics, this cohomology theory upgrades to a six-functor formalism.
\begin{theorem*}[\ref{thm: main six-functor formalism}]
    Let $\kappa$ be an uncountable regular cardinal. There is a six-functor formalism $X\mapsto\QCoh_{\sld}(X)$ on the $(\infty,1)$-category $\Stk_{\kappa}$ of $\kappa$-small spectral algebraic stacks for the faithfully flat and finitely presented topology, assigning to each $X$ its closed symmetric monoidal $(\infty,1)$-category $\QCoh_{\sld}(X)$ of solid quasicoherent sheaves, with the following properties:
    \begin{enumerate}
        \item For every morphism $f: X\to Y$ there is an adjoint pair
        $$f^{*}:\QCoh_{\sld}(Y)\rightleftarrows\QCoh_{\sld}(X): f_{*}$$
        with $f^{*}$ symmetric monoidal.
        \item For every morphism $f: X\to Y$ that locally in the faithfully flat and finitely presented topology factors as an integral morphism followed by a finitely presented morphism, there is a further adjoint pair
        $$f_{!}:\QCoh_{\sld}(X)\rightleftarrows\QCoh_{\sld}(Y): f^{!}$$
        satisfying proper base change and the projection formula.
        \item For every morphism $f: X\to Y$ that locally in the faithfully flat and finitely presented topology factors as $\spec(C)\to\spec(B)\to\spec(A)$ with $A\to B$ flat and almost of finite presentation, and $C$ perfect over $B$, there is a natural equivalence
        $$f^{!}(-)\simeq\omega_{f}\otimes_{X}^{\sld}f^{*}(-):\QCoh_{\sld}(Y)\to\QCoh_{\sld}(X),$$
        where $\omega_{f}:=f^{!}1_{Y}\in\QCoh_{\sld}(X)$ is the dualizing sheaf of $f$.
        \begin{itemize}
            \item[(3.1)] If $f$ is a fiber-smooth map of $\kappa$-small affine spectral schemes, then $\omega_{f}$ is $\otimes_{X}^{\sld}$-invertible.
            \item[(3.2)] If $f$ is \'{e}tale-locally an \'{e}tale morphism, then $\omega_{f}\simeq1_{X}$.
        \end{itemize}
        \item If $f$ is a proper morphism of $\kappa$-small affine spectral schemes, then there is a natural equivalence
        $$f_{!}(-)\simeq f_{*}(-):\QCoh_{\sld}(X)\to\QCoh_{\sld}(Y).$$
        \item If $f: X\to Y$ is as in (3), then the following variant of Grothendieck--Verdier duality holds.
        \begin{itemize}
            \item[(5.1)] The natural morphism $\intHom_{Y}(f_{!}M,N)\to f_{*}\intHom_{X}(M,f^{!}N)$ is an equivalence in $\QCoh_{\sld}(Y)$ for all $M\in\QCoh_{\sld}(X),N\in\QCoh_{\sld}(Y)$.
            \item[(5.2)] The functor given on objects by $M\mapsto\intHom_{X}(M,\omega_{f})$ is an anti-autoequivalence of the full subcategory $\Coh_{\sld}(X,f)$ of $f$-coherent sheaves.
        \end{itemize}
    \end{enumerate}
\end{theorem*}
The theorem sits at the confluence of several areas of mathematics: stable homotopy theory and spectral algebraic geometry, the abstract theory of six-functor formalisms, and condensed mathematics.

In stable homotopy theory the Abelian groups of ordinary algebra give way to the $(\infty,1)$-category $\Sp$ of spectra, the modules over the sphere spectrum $\Sph$, with connective $\EE_{\infty}$-algebras in $\Sp$ taking the place of commutative rings. Following Lurie \cite{SAG}, one builds algebraic geometry over such rings, obtaining spectral analogues of the objects of algebraic geometry -- among them the spectral algebraic stacks and quasicoherent sheaves above.

A sheaf theory is often organized by a six-functor formalism, which records the pushforward and pullback functors, their exceptional counterparts, and the coherences among them. An abstract $(\infty,1)$-categorical account of such formalisms, independent of any one geometric context -- \'{e}tale sheaves, $\Dcal$-modules, and the like -- has emerged in recent work of Gaitsgory--Rozenblyum \cite{GR-DAG}, Heyer--L. Mann \cite{HeyerMann6FF}, Liu--Zheng \cite{LiuZheng6FF}, L. Mann \cite{MannRigid6FF}, and Scholze \cite{Scholze6FF}. One would like the quasicoherent sheaves of spectral algebraic geometry to fit into such a formalism. As in the classical setting, they do not: the exceptional pushforward $f_{!}$ is not defined beyond the case of proper morphisms.

The remedy is due to the condensed mathematics of Clausen and Scholze. They observe that the obstruction can be overcome by enlarging the category of sheaves, replacing quasicoherent sheaves by their topologically complete, or solid, analogues. Within this enlargement, the exceptional pushforward exists for morphisms of finite type \cite[Thm. 8.13]{Condensed} and the six operations assemble into a full six-functor formalism.

The present work carries this out in the setting of spectral algebraic geometry, generalizing the solid six-functor formalism of Clausen and Scholze in the static case \cite{Condensed,Complex} and of Hansen--L. Mann \cite[\S 2.7]{HansenMannCLL}, L. Mann \cite[\S 2.9]{MannRigid6FF}, and Rodr\'{i}guez Camargo \cite{CamargoSolid} in the animated setting. The passage to connective $\EE_{\infty}$ rings is a generalization of the static and animated frameworks, with the forgetful functor from animated rings to connective $\EE_{\infty}$-algebras over $\ZZ$ exhibiting the animated theory as a special case. It is also where the difficulty lies, for the structural conveniences of animated algebra are unavailable over the sphere: the free $\EE_{\infty}$-algebra over the sphere $\Sph\{T\}\simeq\bigoplus_{n\geq0}\Sigma^{\infty}_{+}B\Sigma_{n}$ is of infinite $\Tor$-amplitude, so key factorizations underlying the animated arguments do not transfer.
\subsection{Leitfaden}\label{subsec: intro leitfaden}
We now describe the contents in more detail. The foundation is a theory of analytic rings relative to connective $\EE_{\infty}$-algebras in condensed spectra, generalizing the analytic rings of Clausen--Scholze and their animated refinements, and its construction occupies much of the work. With this theory in place we specialize to the solid setting, building the six operations first on affine spectral schemes and then on stacks, where a form of Grothendieck--Verdier duality is recovered.
\subsubsection{Condensed spectra and solid spectra}\label{subsubsec: intro condensed}
\Cref{sec: condensed spectra and solid spectra} recalls the light condensed mathematics adopted throughout. Unlike in the animated setting, where the theory can be built by applying universal constructions such as animation to the static case, we are forced to work with $(\infty,1)$-categorical constructions from the start. The $(\infty,1)$-categories of condensed animae $\Cond\Ani$ and condensed spectra $\Cond\Sp$ are constructed in \Cref{subsec: condensed animae,subsec: condensed spectra} as hypercomplete sheaves on the site of light profinite sets, and their basic properties are established. Solid spectra are defined in \Cref{subsec: solid spectra}, where the definition of Wagner \cite[\S B.6]{WagnerHHab} through summability of nullsequences is shown to agree with that of Clausen \cite[Def. 13.2]{ClausenLie} through homotopy groups (\Cref{thm: properties of solid spectra}).
\subsubsection{Analytic \texorpdfstring{$\EE_{\infty}$}{E-infinity}-rings}\label{subsubsec: intro analytic}
\Cref{sec: analytic rings} develops the theory of analytic $\EE_{\infty}$-rings. \Cref{subsec: basic definitions} introduces the central object, a condensed $\EE_{\infty}$-ring together with a distinguished subcategory of its modules.
\begin{definition*}[\ref{def: analytic E-infinity ring}]
    Let $A^{\tri}$ be a connective $\EE_{\infty}$-algebra object in $\Cond\Sp$. An \emph{analytic $\EE_{\infty}$-ring} on $A^{\tri}$-modules is a pair $A:=(A^{\tri},\Mod_{A})$ where $\Mod_{A}\subseteq\Mod_{A^{\tri}}$ is a full subcategory of $A^{\tri}$-modules in condensed spectra such that:
    \begin{enumerate}[label=(\roman*)]
        \item $\Mod_{A}$ is closed under small limits and colimits in $\Mod_{A^{\tri}}$.
        \item $\Mod_{A}$ is $\Cond\Sp$-linear: $\intHom_{A^{\tri}}(M,N)\in\Mod_{A}$ for $M\in\Mod_{A^{\tri}},N\in\Mod_{A}$.
        \item The left adjoint $L^{A}:\Mod_{A^{\tri}}\to\Mod_{A}$ to the inclusion is right $t$-exact.
    \end{enumerate}
\end{definition*}
Properties of $A$-analytic modules are studied in \Cref{subsec: analytic modules}, and the functoriality of the $(\infty,1)$-category $\An\CAlg$ of analytic $\EE_{\infty}$-rings, together with its colimits, in \Cref{subsec: functoriality and colimits}. \Cref{subsec: completion} constructs completions (\Cref{thm: completions exist}), the passage to the full subcategory $\An\CAlg^{\wedge}$ of complete analytic $\EE_{\infty}$-rings where $A^{\tri}\in\Mod_{A}\subseteq\Mod_{A^{\tri}}$. This is simpler here than in the animated setting: the completion of an analytic $\EE_{\infty}$-ring is again an analytic $\EE_{\infty}$-ring, whereas endowing the completion of an animated analytic ring with an animated structure is delicate. \Cref{subsec: topological invariance} establishes a topological invariance phenomenon, showing that an analytic $\EE_{\infty}$-ring structure is detected on the heart and so recorded by a purely static datum over $\pi_{0}(A^{\tri})$ (\Cref{thm: heart recognition for analytic rings}). \Cref{subsec: mereology of analytic modules} then turns to the internal finiteness structure of $\Mod_{A}$, relating nuclearity and dualizability (\Cref{prop: dualizable iff nuclear and compact}) as a special case of the discussion in \Cref{app: nuclearity}.
\subsubsection{Morphisms of analytic \texorpdfstring{$\EE_{\infty}$}{E-infinity}-rings}\label{subsubsec: intro morphisms}
\Cref{sec: morphisms of analytic rings} takes up morphisms in earnest. \Cref{subsec: constructing morphisms} records the $\EE_{\infty}$-analogue of a construction principle of Scholze, promoting a morphism of the underlying rings to one of analytic $\EE_{\infty}$-rings under a compatibility on $\pi_{0}$ (\Cref{prop: constructing morphisms of analytic rings}). \Cref{subsec: steady morphisms} treats the steady morphisms.
\begin{definition*}[\ref{def: steady morphism}]
    A morphism $(A^{\tri},\Mod_{A})\to(B^{\tri},\Mod_{B})$ in $\An\CAlg^{\wedge}$ is \emph{steady} if for every morphism $(A^{\tri},\Mod_{A})\to(C^{\tri},\Mod_{C})$ and coCartesian square with pushout $(D^{\tri},\Mod_{D})$ the natural morphism
    $$B\otimes_{A}(M|_{A^{\tri}})\longrightarrow(D\otimes_{C}M)|_{B^{\tri}}$$
    is an equivalence in $\Mod_{B}$ for all $M\in\Mod_{C}$.
\end{definition*}
For a morphism of analytic $\EE_{\infty}$-rings, the scalar extension $\Mod_{A}\to\Mod_{B}$ need not be tensoring with $B^{\tri}$ regarded as an algebra in $\Mod_{A}$, so base change is not automatic. 

\Cref{subsec: open immersions} treats open immersions (\Cref{def: open immersion}) and, in particular, a restricted form of an expectation of Peter Scholze, relayed to us by L. Mann, that open immersions are detected on the heart. The recognition reduces the question to the underlying static analytic rings once the analytic $\Tor$-amplitude is at most one.
\begin{proposition*}[\ref{prop: open immersion detected on static rings}]
    Let $(A^{\tri},\Mod_{A})\to(B^{\tri},\Mod_{B})$ be a morphism of complete analytic $\EE_{\infty}$-rings with $B\otimes_{A}\pi_{0}(A^{\tri})\simeq\pi_{0}(B^{\tri})$ and such that:
    \begin{enumerate}[label=(\roman*)]
        \item The right adjoint $(-)|_{A^{\tri}}:\Mod_{B}\to\Mod_{A}$ to scalar extension is fully faithful.
        \item The functor $\pi_{0}(B)\otimes_{\pi_{0}(A)}-$ is of analytic $\Tor$-amplitude $\leq 1$.
        \item The $t$-structure on $\Mod_{A}$ is compatible with small products, and the induced morphism $(\pi_{0}(A^{\tri}),\Mod_{\pi_{0}(A)})\to(\pi_{0}(B^{\tri}),\Mod_{\pi_{0}(B)})$ is an open immersion.
    \end{enumerate}
    Then $(A^{\tri},\Mod_{A})\to(B^{\tri},\Mod_{B})$ is an open immersion.
\end{proposition*}
\Cref{subsec: analytically proper morphisms} treats the complementary analytically proper morphisms and their base change and projection formulae. Together the open immersions and analytically proper morphisms form a suitable decomposition of a geometric setup on $\An\Aff:=(\An\CAlg^{\wedge})^{\Opp}$ (\Cref{lem: geometric setup on analytic E-infinity rings}), from which \Cref{subsec: six operations on analytic rings} constructs a six-functor formalism on analytic $\EE_{\infty}$-rings (\Cref{prop: six-functor formalism on analytic E-infinity rings}), following the framework of \cite[\S 2-3]{HeyerMann6FF}. \Cref{subsec: descent} closes with a discussion of descent for analytic modules (\Cref{thm: descent for descendable morphisms}), building on work of Mikami \cite[\S 2]{MikamiFFDescent} and L. Mann's formalism of steady endofunctors \cite[\S 2.5-2.6]{MannRigid6FF}.
\subsubsection{Discrete solid \texorpdfstring{$\EE_{\infty}$}{E-infinity}-rings}\label{subsubsec: intro solid}
\Cref{sec: solid commutative rings} begins the study of discrete solid $\EE_{\infty}$-rings. The solid spectra of \Cref{subsec: solid spectra} assemble into an analytic $\EE_{\infty}$-ring over the sphere, the solid sphere $\Sph_{\sld}$ of \Cref{subsec: solid sphere}, whose generators are internally compact projective rather than merely $\omega_{1}$-compact (\Cref{prop: generation of solid spectra}, cf. \Cref{prop: A-analytic tensor product}). The same construction, carried out on compact projective algebras and extended by sifted colimits, produces the discrete solid $\EE_{\infty}$-ring $A_{\sld}$ of an arbitrary connective $\EE_{\infty}$-algebra $A\in\CAlg^{\cn}$ in \Cref{subsec: solid E-infinity rings} (\Cref{prop: solid rings general}). \Cref{subsec: spectral Huber pairs} refines these to discrete spectral Huber pairs.
\begin{definition*}[\ref{def: discrete spectral Huber pair}]
    Let $A\in\CAlg^{\cn}$. A \emph{discrete spectral Huber pair} $(A,A^{+})$ is a morphism $A^{+}\to A$ in $\CAlg^{\cn}$ such that $\pi_{0}(A^{+})\subseteq\pi_{0}(A)$ is an integrally closed subring of $\pi_{0}(A)$. A \emph{morphism of discrete spectral Huber pairs} $f:(A,A^{+})\to(B,B^{+})$ is a morphism $f: A\to B$ in $\CAlg^{\cn}$ with $\pi_{0}(f)(\pi_{0}(A^{+}))\subseteq\pi_{0}(B^{+})$.
\end{definition*}
In contrast to the extant literature on the animated theory, the plus-part $A^{+}$ is permitted to be non-static, the subring condition being imposed only on $\pi_{0}$. Such pairs give rise to complete analytic $\EE_{\infty}$-rings by a fully faithful functor (\Cref{lem: Huber pairs to solid rings,lem: functoriality of Huber rings}), and afford both finer control of the solid structure and a convenient factorization of morphisms through an analytically proper one (\Cref{prop: factorization of Huber ring maps}). \Cref{subsec: mereology of solid modules} shows that over a discrete solid Huber ring the nuclear objects are exactly the discrete ones, so every dualizable object is discrete and $A_{\sld}$ is Fredholm (\Cref{prop: nuclear objects in solid Huber rings,prop: solid rings are Fredholm}).

The exceptional pushforwards rest on the factorization discrete spectral Huber pairs supply. The morphisms carrying them are those of finite expansion (\Cref{def: finite expansion Einfty}), the morphisms $A\to B$ in $\CAlg^{\cn}$ admitting a factorization $A\to A\{T_{1},\dots,T_{n}\}\to B$ through a free $\EE_{\infty}$-algebra with $\pi_{0}(A\{T_{1},\dots,T_{n}\})\to\pi_{0}(B)$ integral. The construction of $f_{!}\dashv f^{!}$ for such a morphism reduces to the free-algebra map $\Sph_{\sld}\to\Sph\{T\}_{\sld}$, treated in \Cref{subsec: solid affine line,subsec: exceptional functors finite expansion}.
\begin{proposition*}[\ref{prop: exceptional functors for the affine line}]
    Let $f:\Sph_{\sld}\to\Sph\{T\}_{\sld}$ be the morphism of discrete solid $\EE_{\infty}$-rings to the free $\EE_{\infty}$-algebra.
    \begin{enumerate}
        \item $f$ admits a factorization $\Sph_{\sld}\xrightarrow{p}(\Sph\{T\},\Sph)_{\sld}\xrightarrow{j}\Sph\{T\}_{\sld}$ in which $p$ is analytically proper and $j$ an open immersion.
        \item The functor $f_{!}:=p_{*}j_{!}:\Mod_{\Sph\{T\}_{\sld}}\to\Mod_{\Sph_{\sld}}$ preserves small colimits and satisfies the projection formula.
    \end{enumerate}
\end{proposition*}
That $j$ is an open immersion is shown in two ways, either through the idempotent algebra of \Cref{subsec: solid affine line} or by the recognition criterion of \Cref{prop: open immersion detected on static rings} described above. Base change gives the free algebra in finitely many variables over an arbitrary ring. Noting that integral morphisms are analytically proper, the two factors yield the exceptional functors of any finite-expansion morphism (\Cref{thm: exceptional functor for finite expansion}). \Cref{subsec: flat morphisms of solid rings} develops the structure theory of the flat, almost of finite presentation morphisms after Mikami, on which \Cref{subsec: descent for solid rings} shows discrete solid $\EE_{\infty}$-rings satisfy descent for the faithfully flat and finitely presented topology.
\begin{theorem*}[\ref{thm: solid fppf descent}]
    Let $A_{\sld}\to B_{\sld}$ be a finitely presented and faithfully flat morphism of discrete solid $\EE_{\infty}$-rings. Then the functor
    $$\Mod_{A_{\sld}}\longrightarrow\lim_{[n]\in\Delta}\Mod_{B_{\tsld}^{\otimes_{A_{\tsld}}(n+1)}}$$
    is an equivalence.
\end{theorem*}
\subsubsection{Six operations for solid quasicoherent sheaves}\label{subsubsec: intro six operations}
\Cref{sec: six functors for solid sheaves} constructs the six-functor formalism for solid quasicoherent sheaves. On $\kappa$-small affine spectral schemes it is obtained in \Cref{subsec: six functors on affines} by restricting the formalism on complete analytic $\EE_{\infty}$-rings along the embedding of discrete solid $\EE_{\infty}$-rings (\Cref{prop: six-functor formalism affine spectral schemes}). 

\Cref{subsec: affine prim and suave} studies the prim, suave, and smooth morphisms of Heyer--L. Mann in this setting, where the divergence from the animated theory is sharpest. In the animated setting, a morphism of solid affine schemes is suave if and only if it is almost of finite presentation and of finite $\Tor$-amplitude. The proof there rests on a structure theorem factoring such a morphism as $A\to A[T_{1},\dots,T_{n}]\to B$ over the polynomial algebra, and no such theorem is available over the sphere. We isolate instead the pseudo-lisse morphisms -- those factoring as $A\to B\to C$ with $A\to B$ flat and almost of finite presentation and $C$ perfect over $B$ -- and prove they are $\msld$-suave (\Cref{prop: affine suave morphisms}), though it is not clear that all $\msld$-suave morphisms are of this form (\Cref{rmk: comparison to lurie jiang}). 

\Cref{subsec: affine perfect and coherent} introduces the perfect and $f$-coherent sheaves and their duality. \Cref{subsec: construction on stacks} globalizes the affine formalism to spectral algebraic stacks by descent for the faithfully flat and finitely presented topology (\Cref{thm: six-functor formalism on stacks}), using that such morphisms are $\msld$-suave (\Cref{lem: ffp is universal cover}) to reduce many questions to the affine case. \Cref{subsec: duality theory for spectral algebraic stacks} then develops Grothendieck--Verdier duality (\Cref{thm: Grothendieck Verdier duality stacks}) and restates the main theorem.
\subsubsection{Appendix A: \texorpdfstring{$t$}{t}-bistable categories and their localizations}\label{subsubsec: intro appendix A}
\Cref{app: categories} collects results in $(\infty,1)$-category theory that we could not find in the literature. Several statements usually proved for analytic rings already hold for stable presentable categories carrying compatible $t$- and symmetric monoidal structures -- the $t$-bistable categories. We work at this level of generality, both to clarify the hypotheses on which the statements rest, and because we expect them to be of independent interest. The localizations of such categories, and the recognition of symmetric monoidal localizations with smashing kernel, imply many of the topological invariance results of \Cref{subsec: topological invariance} and underlie the previously described open-immersion criterion of \Cref{subsec: open immersions}.
\subsubsection{Appendix B: Nuclearity in the \texorpdfstring{$\kappa$}{kappa}-presentable setting}\label{subsubsec: intro appendix B}
\Cref{app: nuclearity} treats trace class morphisms and nuclear objects in stable $\kappa$-presentably symmetric monoidal $(\infty,1)$-categories with compact unit, showing that the compact-generation hypotheses of \cite[Lectures VIII \& IX]{Complex} are unnecessary for the categories of interest. The nuclear objects (\Cref{def: basic nuclear in category}) form a full subcategory closed under small colimits, and an object is dualizable exactly when it is compact and nuclear (\Cref{prop: dualizable iff nuclear and compact category}).
\subsubsection{Appendix C: Commutative algebra of \texorpdfstring{$\EE_{\infty}$}{E-infinity}-rings}\label{subsubsec: intro appendix C}
\Cref{app: commutative algebra} records auxiliary results in the commutative algebra of connective $\EE_{\infty}$-algebras in $\Sp$. Finiteness properties of modules and morphisms are recalled in \Cref{subsec: finiteness properties}. We then introduce in \Cref{subsec: morphisms of finite expansion} the morphisms of finite expansion, due to Hamacher \cite{FiniteExpansion} in the static case. We adapt this to the $\EE_{\infty}$ setting and verify its permanence properties (\Cref{prop: finite expansion permanence}). \Cref{subsec: pseudo-lisse morphisms} concludes with the pseudo-lisse morphisms $A\to C$ of connective $\EE_{\infty}$-algebras in $\Sp$, those admitting a factorization $A\to B\to C$ with $B$ flat and almost of finite presentation over $A$ and $C$ perfect as a $B$-module, alongside the elementary pseudo-lisse morphisms, where $B$ is required to be the polynomial algebra over $A$ in finitely many variables (\Cref{def: pseudo-lisse appendix}), and their basic properties (\Cref{prop: properties of pseudo-lisse morphisms}).
\subsection{Notation and conventions}\label{subsec: notation and conventions}
In the rest of this text, we will work with $\infty$-categories and $\infty$-categorical algebra throughout following \cite{HTT,HA,kerodon}, referring to $(\infty,1)$-categories as categories and distinguishing ordinary categories as needed. We will adopt the following: 
\begin{itemize}
  \item $\Ani$ and $\Sp$ will refer to the categories of animae (also known as spaces) and spectra, respectively. 
  \item $\PrL$ will refer to Lurie's category of presentable categories and left adjoint functors. For a regular cardinal $\kappa$, $\PrL_{\kappa}$ will refer to Lurie's category of $\kappa$-presentable categories -- the non-full subcategory of $\PrL$ spanned by the $\kappa$-presentable categories and functors preserving $\kappa$-compact objects. $\PrL$ admits a symmetric monoidal structure via the Lurie tensor product which restricts to a symmetric monoidal structure on $\PrL_{\kappa}$.
  \item We say a category $\Cscr$ is presentably symmetric monoidal if $\Cscr$ is symmetric monoidal and the tensor product preserves small colimits separately in each argument. This implies that $\Cscr$ is closed symmetric monoidal. For a regular cardinal $\kappa$, we say $\Cscr$ is $\kappa$-presentably symmetric monoidal if $\Cscr\in\CAlg(\PrL_{\kappa})$: that is, the unit $1_{\Cscr}$ is $\kappa$-compact and the tensor product preserves $\kappa$-compact objects. 
\end{itemize}
Additionally, we will make use of the theory of (light) condensed mathematics as developed by Clausen and Scholze, the language of enriched categories (ie. enriched $\infty$-categories), as well as the axiomatic perspective on six-functor formalisms. We refer the reader to \cite{Condensed,Analytic,Complex,AnalyticStacks,CamargoSolid} as references for (light) condensed mathematics, \cite{MannRigid6FF,HeyerMann6FF,RamziLocRigid,RamziDualizable,ReutterZettoEnriched} for enriched categories, and \cite{MannRigid6FF,HeyerMann6FF,ConstructionCLL} for abstract six-functor formalisms. We will adopt the following: 
\begin{itemize}
  \item Condensed objects are assumed to be light, that is, hypersheaves on the category $\PFin_{\NN}$ of profinite sets arising as sequential limits of finite sets and continuous maps between them.
  \item An object is said to be static if it is 0-truncated as an object (eg. if it lies in the heart of a $t$-structured stable category), reserving discreteness for condensed objects satisfying a certain property. 
  \item For a category $\Cscr$ enriched in a category $\Vscr$, we will denote by $\Mor_{\Cscr}(-,-)$ and $\Hom_{\Cscr}^{\Vscr}(-,-)$ the $\Ani$-valued and $\Vscr$-valued morphisms, respectively. When $\Cscr$ is stable and $\Vscr=\Sp$, we write $\Hom_{\Cscr}^{\Sp}(-,-)$ as $\Hom_{\Cscr}(-,-)$. If $\Cscr$ is closed symmetric monoidal (eg. presentably symmetric monoidal) we will use $\intHom_{\Cscr}(-,-)$ to refer to the mapping object internal to $\Cscr$.
\end{itemize}
In this text, we have at times elected to deviate our nomenclature and notation from the extant literature. We have done so only where, in our judgment, the alternative conventions render the underlying ideas more transparent -- first to ourselves, and, we hope, to the reader. In each such instance, the justification for our choice is given in the remark immediately succeeding the definition in which the deviation is introduced. We apologize in advance for any unnecessary confusion these departures may cause the reader already accustomed to the established terminology.
\subsection{Acknowledgements}\label{subsec: acknowledgements}
I would like to thank my advisors: Dr. Lucas Mann, for his most felicitous choice of topic for my master's thesis, for his generosity in discussion and critique, and for his unwavering support throughout the year; and Prof. Dr. Peter Scholze, for the inspiring courses he taught as well as for helpful conversations and correspondence. I have been fortunate to discuss the content of this thesis, in part or in whole, with various mathematicians. I thank Benjamin Antieau, Daniel Arone, Sam Brinkerhoff, Dustin Clausen, Thorger Gei{\ss}, Niklas Kipp, Jordan Levin, Jiacheng Liang, Niko Naumann, Marius Nielsen, Lucas Piessevaux, Phil P\"{u}tzst\"{u}ck, Maxime Ramzi, Gabriel Ribeiro, Florian Riedel, Maria Stroe, Ferdinand Wagner, and Yuchen Wu for conversations that have shaped this work, even if not always in its content. I am particularly grateful to Maria Stroe, who read a draft with great care.